\begin{table}[t]
\caption{Distractor Scaling: Questions Fixed, Pool Grown}
\label{tab:scaling}
\centering
\footnotesize
\setlength{\tabcolsep}{3.5pt}
\begin{tabular}{@{}llcccccc@{}}
\toprule
\textbf{Pool} & \textbf{System} & \textbf{Hit@5} & \textbf{MRR} & \textbf{Evid.} & \textbf{Stale@1} & \textbf{Dup.} & \textbf{Leak} \\
\midrule
15 mem. & Hybrid & 1.000 & 0.944 & 1.000 & 0.182 & 0.022 & 0.022 \\
 & Graph-aug. & 1.000 & 0.944 & 1.000 & 0.182 & 0.022 & 0.022 \\
 & Adaptive & 1.000 & 1.000 & 1.000 & 0.000 & 0.000 & 0.000 \\
40 mem. & Hybrid & 1.000 & 0.944 & 1.000 & 0.182 & 0.170 & 0.021 \\
 & Graph-aug. & 1.000 & 0.944 & 1.000 & 0.182 & 0.192 & 0.021 \\
 & Adaptive & 1.000 & 1.000 & 1.000 & 0.000 & 0.000 & 0.000 \\
65 mem. & Hybrid & 1.000 & 0.944 & 1.000 & 0.182 & 0.192 & 0.021 \\
 & Graph-aug. & 1.000 & 0.944 & 1.000 & 0.182 & 0.213 & 0.021 \\
 & Adaptive & 1.000 & 1.000 & 1.000 & 0.000 & 0.000 & 0.000 \\
115 mem. & Hybrid & 1.000 & 0.889 & 1.000 & 0.273 & 0.234 & 0.021 \\
 & Graph-aug. & 1.000 & 0.889 & 1.000 & 0.273 & 0.255 & 0.021 \\
 & Adaptive & 1.000 & 1.000 & 1.000 & 0.000 & 0.000 & 0.000 \\
\bottomrule
\end{tabular}
\par\vspace{2pt}
\raggedright\scriptsize Evidence recall is the fraction of required gold memories returned for multi-evidence questions. Distractors are drawn from the same vocabulary as the corpus but contain no gold key, so they can only cost a system rank.
\end{table}
